\section{Execution Cycle \& Memory Handling}
The interaction between the deterministic orchestration layer and the generative model is governed by a stateful execution cycle. Unlike standard single-turn LLM interactions, the Galois framework maintains a persistence layer across iterations to maximize recall and prevent redundancy. This logic is encapsulated within the \textbf{GaloisExecutor} class and its specialized internal memory component, \textbf{GaloisMemory}.

\subsection{The $\texttt{GaloisMemory}$ Management Class}

\subsubsection{Purpose and Architecture}
The $\texttt{GaloisMemory}$ class serves as the dedicated state management component for the physical execution strategies ($\texttt{Table-Scan}$ and $\texttt{Key-Scan}$). Inheriting from $\texttt{BaseMemory}$ and $\texttt{BaseModel}$, it is architected to fulfill three primary responsibilities: maintaining an accumulated history of unique records and their logprobs extracted across multiple iterative steps, enforcing strict deduplication to preserve data integrity, and formatting these accumulated records into a structured context ($\mathcal{H}$) for injection into subsequent LLM prompts. By managing this persistent state, the class ensures that the model is continuously guided toward novel data discovery while preventing redundant processing.

\subsubsection{Internal State Variables}
The memory state is maintained by a set of core attributes defined as Pydantic fields. These variables track the data accumulation, deduplication status, and execution metrics:

% Richiede \usepackage{booktabs} nel preambolo
\begin{table}[H]
    \centering
    \renewcommand{\arraystretch}{1.2}
    \begin{tabular}{p{0.25\linewidth} p{0.70\linewidth}}
        \toprule
        \textbf{Attribute \& Type} & \textbf{Description} \\
        \midrule
        \textbf{memory}\newline \small{\texttt{List[Dict[str, Any]]}} & 
        The primary storage structure containing the complete collection of unique records (tuples) retrieved thus far. \\
        \addlinespace
        \textbf{key\_values}\newline \small{\texttt{set}} & 
        A hash set storing unique key value combinations (as Python tuples). \\
        \addlinespace
        \textbf{tokens}\newline \small{\texttt{int}} & 
        An aggregate counter tracking the total number of tokens processed by the LLM across all invocations. \\
        \addlinespace
        \textbf{time}\newline \small{\texttt{float}} & 
        An aggregate measure of the total processing time (latency) associated with the LLM interactions. \\
        \addlinespace
        \textbf{logprobs}\newline \small{\texttt{List[float]}} & 
        A list storing the average log-probabilities for each generated record, used to estimate confidence. \\
        \bottomrule
    \end{tabular}
    \caption{Internal State Attributes of GaloisMemory}
    \label{tab:memory_attributes}
\end{table}

\subsubsection{Core Methods and Functionality}
The functional logic of the memory component is implemented through four distinct methods that handle the lifecycle of data persistence.
\\\\
The \textbf{save\_context(inputs, outputs)} method acts as the central mechanism for processing new records generated by an LLM invocation. It iterates through the provided response records and uses the internal helper \textbf{\_get\_key\_values(record, keys)} to extract composite key tuples. These keys are checked against the \texttt{self.key\_values} set; if a key is absent, the record is appended to \texttt{self.memory} and the key is registered.\\ 
Simultaneously, the method preserves the pre-calculated average log-probabilities associated with each tuple into the \texttt{self.logprobs} structure, maintaining a strict one-to-one correspondence between the stored records and their confidence scores for subsequent quality filtering.
\\\\
Crucially, this method implements the convergence logic: if the size of the key set remains unchanged after processing, it raises a $\texttt{NoNewTuplesFound}$ exception to signal the termination of the iterative process. Additionally, it aggregates execution metrics ($\texttt{tokens}$ and $\texttt{time}$) from the output.
\\\\
To support the iterative prompting cycle, the $\textbf{load\_memory\_variables(inputs)}$ method prepares the accumulated records for injection into the next prompt. It serializes the current $\texttt{self.memory}$ list into a compact JSON string, returning it under the key $\texttt{"history"}$. 
\\\\
Finally, the $\textbf{clear()}$ method provides a reset mechanism, emptying the storage structures and zeroing out metrics to ensure a clean state for subsequent query executions.


\subsection{The Iterative Execution Loop and Termination Criteria}

Both the \texttt{Table-Scan} and \texttt{Key-Scan} physical operators rely on a fundamental Iterative Execution Loop (lines 4-19 in Algorithms \ref{alg:table_scan} and \ref{alg:key_scan}) to overcome the LLM context window limitations and maximize data extraction coverage. The core mechanism involves repeatedly querying the LLM and dynamically updating the prompt with historical data.

\subsubsection{Loop Mechanism: Pagination \& History Injection}
In both strategies, the iterative process is driven by the dynamic insertion of previously acquired data into the subsequent prompt. This mechanism ensures two critical functions: \textbf{Deduplication}, where the LLM is implicitly guided to avoid regenerating records already present in the injected history ($\mathcal{H}$), and \textbf{Novelty Search}, where the context directs the LLM to explore uncollected portions of the solution space to maximize the cardinality of the final result set ($\mathcal{R}_{final}$).
\\\\
To enhance retrieval efficiency and prevent the LLM from wandering randomly, the iterative mechanism implements a \textbf{value-based pagination} strategy. In addition to the history, the system identifies the last retrieved value (\texttt{last\_val}) from the memory. The prompt logic explicitly instructs the LLM to generate unique tuples that strictly follow \texttt{last\_val} (alphabetically or numerically). This transforms the extraction into a pseudo-sequential scan, significantly improving coverage compared to simple history injection. The primary distinction between the operators lies in the content of the history: \texttt{Table-Scan} injects a history of complete tuples, while \texttt{Key-Scan} injects only the unique key values ($\mathcal{K}$) collected so far.

\subsection{Termination Criteria (Convergence Conditions)}
The iterative loop is governed by two distinct termination criteria, ensuring both safety against infinite execution and functional completeness. The system employs a \texttt{break} statement within the execution loop immediately when either condition is met.
\\\\
The first criterion is the \textbf{Maximum Iteration Limit ($I_{max}$)}. This condition is triggered when the loop counter ($i$) reaches a maximum permissible value, which is configurable via the system's YAML parameter file. The primary purpose of this limit is to serve as a resource fail-safe, strictly bounding token consumption and total execution latency. It prevents the system from entering infinite loops or wasting computational resources in scenarios where the LLM might continue to generate redundant or low-quality data indefinitely.
\\\\
The second criterion is \textbf{Data Convergence}, which serves as the primary functional exit condition. This state is reached when the system determines that an LLM invocation has yielded \textbf{no new unique records} (for \texttt{Table-Scan}) or \textbf{no new unique keys} (for \texttt{Key-Scan}). This condition is explicitly detected by the \texttt{save\_context} function within the \texttt{GaloisMemory} class, which monitors the growth of the internal key set. When the set size remains static after a response, it signals that the LLM has likely exhausted its capacity to discover novel data relevant to the query, allowing the loop to terminate prematurely to maximize efficiency.
\\\\
The system employs a \texttt{break} statement within the execution loop immediately when convergence is detected, effectively overriding the $I_{max}$ constraint if full coverage is achieved earlier.